% Pipeline figure for Music-CRS paper
% Compile standalone or include in main.tex with % Pipeline figure for Music-CRS paper
% Compile standalone or include in main.tex with % Pipeline figure for Music-CRS paper
% Compile standalone or include in main.tex with % Pipeline figure for Music-CRS paper
% Compile standalone or include in main.tex with \input{fig_pipeline.tex}
% Requires: \usepackage{tikz} \usetikzlibrary{positioning, fit, backgrounds}
\begin{figure*}[t]
\centering
\resizebox{\textwidth}{!}{%
\begin{tikzpicture}[
    node distance=0.4cm and 0.3cm,
    box/.style={draw, rounded corners=3pt, align=center, minimum height=0.9cm, font=\small},
    stage/.style={draw, rounded corners=5pt, dashed, inner sep=8pt, thick, fill=gray!5},
    leg/.style={box, fill=blue!10, minimum width=1.8cm, font=\scriptsize},
    note/.style={font=\scriptsize, text=gray, anchor=west},
    arrow/.style={->, thick, >=stealth},
]

% Input
\node[box, fill=yellow!15, minimum width=12cm, minimum height=1.2cm] (input) {
    \begin{tabular}{c}
    \textbf{Multi-Turn Conversation} \\[-2pt]
    {\scriptsize User: ``Play something chill from the 90s'' $\to$ Rec $\to$ User: ``More like that but slower''}
    \end{tabular}
};

% Stage 1: retrieval legs
\node[leg, below=1.4cm of input, xshift=-5cm] (bm25) {BM25\\(sparse)};
\node[leg, right=0.2cm of bm25] (dense) {Dense Emb\\(meta/cf/img)};
\node[leg, right=0.2cm of dense] (pmi) {PMI\\co-occur.};
\node[leg, right=0.2cm of pmi] (bienc) {BGE / E5 /\\Stella ($\times$5)};
\node[leg, right=0.2cm of bienc] (nvemb) {NV-Embed\\v2 (7.8B)};
\node[leg, right=0.2cm of nvemb] (mq) {Multi-\\query ($\times$2)};

% Stage 1 title sits INSIDE the dashed box, so the incoming arrow (which stops
% at the box border) cannot run through it.
\node[font=\small\bfseries, anchor=south west] (s1title)
    at ([yshift=4pt]bm25.north west) {Stage 1: Diverse Retrieval (13 legs $\times$ top-100)};

\begin{scope}[on background layer]
\node[stage, fit=(s1title)(bm25)(dense)(pmi)(bienc)(nvemb)(mq)] (s1box) {};
\end{scope}

% Arrow from input to stage 1
\draw[arrow] (input.south) -- (s1box.north);

% Union node
\node[box, below=0.5cm of s1box, fill=green!10, minimum width=5cm] (union) {
    Union $\to$ \textbf{$\sim$300 candidates/turn} (gold-in-pool $\approx$ 60\%)
};
\draw[arrow] (s1box.south) -- (union.north);

% Stage 2
\node[box, below=0.5cm of union, fill=orange!12, minimum width=9cm, minimum height=1.4cm] (lgbm) {
    \begin{tabular}{c}
    \textbf{Stage 2: LightGBM LambdaRank} \\[-2pt]
    {\scriptsize Features: rank \& cosine score per leg, popularity, turn\#, PMI sum, \#legs containing}\\[-2pt]
    {\scriptsize $\to$ Output: Top-20 ranked tracks}
    \end{tabular}
};
\draw[arrow] (union.south) -- (lgbm.north);

% Stage 3
\node[box, below=0.5cm of lgbm, fill=purple!10, minimum width=7cm, minimum height=1cm] (resp) {
    \begin{tabular}{c}
    \textbf{Stage 3: Qwen3-0.6B Response Generation} \\[-2pt]
    {\scriptsize Track metadata $\to$ grounded natural language explanation}
    \end{tabular}
};
\draw[arrow] (lgbm.south) -- (resp.north);

% Output
\node[box, below=0.4cm of resp, fill=yellow!15, minimum width=6cm] (output) {
    \textbf{Submission}: 20 ranked tracks + response text
};
\draw[arrow] (resp.south) -- (output.north);

% Annotations, left-aligned in a single column just outside the widest box
\coordinate (ann) at ([xshift=0.35cm]s1box.east);
\node[note] at (ann)          {Hit@100$\approx$48\%};
\node[note] at (ann |- union) {Hit@20=39.1\%};
\node[note] at (ann |- lgbm)  {nDCG@20=0.185};

\end{tikzpicture}
}
\caption{System architecture. Stage 1 employs 13 diverse retrieval legs with different model architectures and query formulations. Their union forms a candidate pool with $\sim$60\% gold-track coverage. Stage 2 applies LightGBM LambdaRank to re-rank candidates using per-leg features. Stage 3 generates metadata-grounded responses.}
\label{fig:pipeline}
\end{figure*}

% Requires: \usepackage{tikz} \usetikzlibrary{positioning, fit, backgrounds}
\begin{figure*}[t]
\centering
\resizebox{\textwidth}{!}{%
\begin{tikzpicture}[
    node distance=0.4cm and 0.3cm,
    box/.style={draw, rounded corners=3pt, align=center, minimum height=0.9cm, font=\small},
    stage/.style={draw, rounded corners=5pt, dashed, inner sep=8pt, thick, fill=gray!5},
    leg/.style={box, fill=blue!10, minimum width=1.8cm, font=\scriptsize},
    note/.style={font=\scriptsize, text=gray, anchor=west},
    arrow/.style={->, thick, >=stealth},
]

% Input
\node[box, fill=yellow!15, minimum width=12cm, minimum height=1.2cm] (input) {
    \begin{tabular}{c}
    \textbf{Multi-Turn Conversation} \\[-2pt]
    {\scriptsize User: ``Play something chill from the 90s'' $\to$ Rec $\to$ User: ``More like that but slower''}
    \end{tabular}
};

% Stage 1: retrieval legs
\node[leg, below=1.4cm of input, xshift=-5cm] (bm25) {BM25\\(sparse)};
\node[leg, right=0.2cm of bm25] (dense) {Dense Emb\\(meta/cf/img)};
\node[leg, right=0.2cm of dense] (pmi) {PMI\\co-occur.};
\node[leg, right=0.2cm of pmi] (bienc) {BGE / E5 /\\Stella ($\times$5)};
\node[leg, right=0.2cm of bienc] (nvemb) {NV-Embed\\v2 (7.8B)};
\node[leg, right=0.2cm of nvemb] (mq) {Multi-\\query ($\times$2)};

% Stage 1 title sits INSIDE the dashed box, so the incoming arrow (which stops
% at the box border) cannot run through it.
\node[font=\small\bfseries, anchor=south west] (s1title)
    at ([yshift=4pt]bm25.north west) {Stage 1: Diverse Retrieval (13 legs $\times$ top-100)};

\begin{scope}[on background layer]
\node[stage, fit=(s1title)(bm25)(dense)(pmi)(bienc)(nvemb)(mq)] (s1box) {};
\end{scope}

% Arrow from input to stage 1
\draw[arrow] (input.south) -- (s1box.north);

% Union node
\node[box, below=0.5cm of s1box, fill=green!10, minimum width=5cm] (union) {
    Union $\to$ \textbf{$\sim$300 candidates/turn} (gold-in-pool $\approx$ 60\%)
};
\draw[arrow] (s1box.south) -- (union.north);

% Stage 2
\node[box, below=0.5cm of union, fill=orange!12, minimum width=9cm, minimum height=1.4cm] (lgbm) {
    \begin{tabular}{c}
    \textbf{Stage 2: LightGBM LambdaRank} \\[-2pt]
    {\scriptsize Features: rank \& cosine score per leg, popularity, turn\#, PMI sum, \#legs containing}\\[-2pt]
    {\scriptsize $\to$ Output: Top-20 ranked tracks}
    \end{tabular}
};
\draw[arrow] (union.south) -- (lgbm.north);

% Stage 3
\node[box, below=0.5cm of lgbm, fill=purple!10, minimum width=7cm, minimum height=1cm] (resp) {
    \begin{tabular}{c}
    \textbf{Stage 3: Qwen3-0.6B Response Generation} \\[-2pt]
    {\scriptsize Track metadata $\to$ grounded natural language explanation}
    \end{tabular}
};
\draw[arrow] (lgbm.south) -- (resp.north);

% Output
\node[box, below=0.4cm of resp, fill=yellow!15, minimum width=6cm] (output) {
    \textbf{Submission}: 20 ranked tracks + response text
};
\draw[arrow] (resp.south) -- (output.north);

% Annotations, left-aligned in a single column just outside the widest box
\coordinate (ann) at ([xshift=0.35cm]s1box.east);
\node[note] at (ann)          {Hit@100$\approx$48\%};
\node[note] at (ann |- union) {Hit@20=39.1\%};
\node[note] at (ann |- lgbm)  {nDCG@20=0.185};

\end{tikzpicture}
}
\caption{System architecture. Stage 1 employs 13 diverse retrieval legs with different model architectures and query formulations. Their union forms a candidate pool with $\sim$60\% gold-track coverage. Stage 2 applies LightGBM LambdaRank to re-rank candidates using per-leg features. Stage 3 generates metadata-grounded responses.}
\label{fig:pipeline}
\end{figure*}

% Requires: \usepackage{tikz} \usetikzlibrary{positioning, fit, backgrounds}
\begin{figure*}[t]
\centering
\resizebox{\textwidth}{!}{%
\begin{tikzpicture}[
    node distance=0.4cm and 0.3cm,
    box/.style={draw, rounded corners=3pt, align=center, minimum height=0.9cm, font=\small},
    stage/.style={draw, rounded corners=5pt, dashed, inner sep=8pt, thick, fill=gray!5},
    leg/.style={box, fill=blue!10, minimum width=1.8cm, font=\scriptsize},
    note/.style={font=\scriptsize, text=gray, anchor=west},
    arrow/.style={->, thick, >=stealth},
]

% Input
\node[box, fill=yellow!15, minimum width=12cm, minimum height=1.2cm] (input) {
    \begin{tabular}{c}
    \textbf{Multi-Turn Conversation} \\[-2pt]
    {\scriptsize User: ``Play something chill from the 90s'' $\to$ Rec $\to$ User: ``More like that but slower''}
    \end{tabular}
};

% Stage 1: retrieval legs
\node[leg, below=1.4cm of input, xshift=-5cm] (bm25) {BM25\\(sparse)};
\node[leg, right=0.2cm of bm25] (dense) {Dense Emb\\(meta/cf/img)};
\node[leg, right=0.2cm of dense] (pmi) {PMI\\co-occur.};
\node[leg, right=0.2cm of pmi] (bienc) {BGE / E5 /\\Stella ($\times$5)};
\node[leg, right=0.2cm of bienc] (nvemb) {NV-Embed\\v2 (7.8B)};
\node[leg, right=0.2cm of nvemb] (mq) {Multi-\\query ($\times$2)};

% Stage 1 title sits INSIDE the dashed box, so the incoming arrow (which stops
% at the box border) cannot run through it.
\node[font=\small\bfseries, anchor=south west] (s1title)
    at ([yshift=4pt]bm25.north west) {Stage 1: Diverse Retrieval (13 legs $\times$ top-100)};

\begin{scope}[on background layer]
\node[stage, fit=(s1title)(bm25)(dense)(pmi)(bienc)(nvemb)(mq)] (s1box) {};
\end{scope}

% Arrow from input to stage 1
\draw[arrow] (input.south) -- (s1box.north);

% Union node
\node[box, below=0.5cm of s1box, fill=green!10, minimum width=5cm] (union) {
    Union $\to$ \textbf{$\sim$300 candidates/turn} (gold-in-pool $\approx$ 60\%)
};
\draw[arrow] (s1box.south) -- (union.north);

% Stage 2
\node[box, below=0.5cm of union, fill=orange!12, minimum width=9cm, minimum height=1.4cm] (lgbm) {
    \begin{tabular}{c}
    \textbf{Stage 2: LightGBM LambdaRank} \\[-2pt]
    {\scriptsize Features: rank \& cosine score per leg, popularity, turn\#, PMI sum, \#legs containing}\\[-2pt]
    {\scriptsize $\to$ Output: Top-20 ranked tracks}
    \end{tabular}
};
\draw[arrow] (union.south) -- (lgbm.north);

% Stage 3
\node[box, below=0.5cm of lgbm, fill=purple!10, minimum width=7cm, minimum height=1cm] (resp) {
    \begin{tabular}{c}
    \textbf{Stage 3: Qwen3-0.6B Response Generation} \\[-2pt]
    {\scriptsize Track metadata $\to$ grounded natural language explanation}
    \end{tabular}
};
\draw[arrow] (lgbm.south) -- (resp.north);

% Output
\node[box, below=0.4cm of resp, fill=yellow!15, minimum width=6cm] (output) {
    \textbf{Submission}: 20 ranked tracks + response text
};
\draw[arrow] (resp.south) -- (output.north);

% Annotations, left-aligned in a single column just outside the widest box
\coordinate (ann) at ([xshift=0.35cm]s1box.east);
\node[note] at (ann)          {Hit@100$\approx$48\%};
\node[note] at (ann |- union) {Hit@20=39.1\%};
\node[note] at (ann |- lgbm)  {nDCG@20=0.185};

\end{tikzpicture}
}
\caption{System architecture. Stage 1 employs 13 diverse retrieval legs with different model architectures and query formulations. Their union forms a candidate pool with $\sim$60\% gold-track coverage. Stage 2 applies LightGBM LambdaRank to re-rank candidates using per-leg features. Stage 3 generates metadata-grounded responses.}
\label{fig:pipeline}
\end{figure*}

% Requires: \usepackage{tikz} \usetikzlibrary{positioning, fit, backgrounds}
\begin{figure*}[t]
\centering
\resizebox{\textwidth}{!}{%
\begin{tikzpicture}[
    node distance=0.4cm and 0.3cm,
    box/.style={draw, rounded corners=3pt, align=center, minimum height=0.9cm, font=\small},
    stage/.style={draw, rounded corners=5pt, dashed, inner sep=8pt, thick, fill=gray!5},
    leg/.style={box, fill=blue!10, minimum width=1.8cm, font=\scriptsize},
    note/.style={font=\scriptsize, text=gray, anchor=west},
    arrow/.style={->, thick, >=stealth},
]

% Input
\node[box, fill=yellow!15, minimum width=12cm, minimum height=1.2cm] (input) {
    \begin{tabular}{c}
    \textbf{Multi-Turn Conversation} \\[-2pt]
    {\scriptsize User: ``Play something chill from the 90s'' $\to$ Rec $\to$ User: ``More like that but slower''}
    \end{tabular}
};

% Stage 1: retrieval legs
\node[leg, below=1.4cm of input, xshift=-5cm] (bm25) {BM25\\(sparse)};
\node[leg, right=0.2cm of bm25] (dense) {Dense Emb\\(meta/cf/img)};
\node[leg, right=0.2cm of dense] (pmi) {PMI\\co-occur.};
\node[leg, right=0.2cm of pmi] (bienc) {BGE / E5 /\\Stella ($\times$5)};
\node[leg, right=0.2cm of bienc] (nvemb) {NV-Embed\\v2 (7.8B)};
\node[leg, right=0.2cm of nvemb] (mq) {Multi-\\query ($\times$2)};

% Stage 1 title sits INSIDE the dashed box, so the incoming arrow (which stops
% at the box border) cannot run through it.
\node[font=\small\bfseries, anchor=south west] (s1title)
    at ([yshift=4pt]bm25.north west) {Stage 1: Diverse Retrieval (13 legs $\times$ top-100)};

\begin{scope}[on background layer]
\node[stage, fit=(s1title)(bm25)(dense)(pmi)(bienc)(nvemb)(mq)] (s1box) {};
\end{scope}

% Arrow from input to stage 1
\draw[arrow] (input.south) -- (s1box.north);

% Union node
\node[box, below=0.5cm of s1box, fill=green!10, minimum width=5cm] (union) {
    Union $\to$ \textbf{$\sim$300 candidates/turn} (gold-in-pool $\approx$ 60\%)
};
\draw[arrow] (s1box.south) -- (union.north);

% Stage 2
\node[box, below=0.5cm of union, fill=orange!12, minimum width=9cm, minimum height=1.4cm] (lgbm) {
    \begin{tabular}{c}
    \textbf{Stage 2: LightGBM LambdaRank} \\[-2pt]
    {\scriptsize Features: rank \& cosine score per leg, popularity, turn\#, PMI sum, \#legs containing}\\[-2pt]
    {\scriptsize $\to$ Output: Top-20 ranked tracks}
    \end{tabular}
};
\draw[arrow] (union.south) -- (lgbm.north);

% Stage 3
\node[box, below=0.5cm of lgbm, fill=purple!10, minimum width=7cm, minimum height=1cm] (resp) {
    \begin{tabular}{c}
    \textbf{Stage 3: Qwen3-0.6B Response Generation} \\[-2pt]
    {\scriptsize Track metadata $\to$ grounded natural language explanation}
    \end{tabular}
};
\draw[arrow] (lgbm.south) -- (resp.north);

% Output
\node[box, below=0.4cm of resp, fill=yellow!15, minimum width=6cm] (output) {
    \textbf{Submission}: 20 ranked tracks + response text
};
\draw[arrow] (resp.south) -- (output.north);

% Annotations, left-aligned in a single column just outside the widest box
\coordinate (ann) at ([xshift=0.35cm]s1box.east);
\node[note] at (ann)          {Hit@100$\approx$48\%};
\node[note] at (ann |- union) {Hit@20=39.1\%};
\node[note] at (ann |- lgbm)  {nDCG@20=0.185};

\end{tikzpicture}
}
\caption{System architecture. Stage 1 employs 13 diverse retrieval legs with different model architectures and query formulations. Their union forms a candidate pool with $\sim$60\% gold-track coverage. Stage 2 applies LightGBM LambdaRank to re-rank candidates using per-leg features. Stage 3 generates metadata-grounded responses.}
\label{fig:pipeline}
\end{figure*}
